\section{Introduction}
\begin{comment}
Understanding critical combustion events is a complex task that involves extensive analysis of both experimental and numerical data. Finding patterns in these datasets is an important step towards building better models for combustin phoenomena. 
 Many tools exist both in the reactive and non reactive flows. Proper Orthogonal Decomposition (POD) (CITAZIONE), is an Singular Value Decomposition (SVD) based method that provide a low rank rappresentation of the flow rewriting it as a sum of orthogonal spatial patterns.
 Principal component analysis has been successfully employed to project high-dimensional data onto a lower-dimensional orthogonal subspace of thermodynamical variables (CITAZIONE), effectively reducing the complexity of the system while retaining its essential features. However, this method is not able to capture the temporal evolution of the system and it is not able to provide a physical interpretation of the modes.
 Dynamic Mode Decomposition (DMD) (CITAZIONE) allows the reduction of the system in spatio-temporal modes which have both a spatial structure and an associated temporal frequency and growth rate. Thus the system is rapresented as the sum of modes evolving with a linear dynamics in time that allows the study of the spectral properties of the sytesm. Although powerful this method suffers from noisy data and it oftern is inadequate to provide an accurate description of turbulent reactive flows. High Order Dynamic Mode Decomposition (HODMD) (CITAZIONE) is a extention of DMD that improves the method by introducing the higher order Koopman hypothesis. This makes it more robust and has already been succesfully applied in reactive flows (CITAZIONE)
Combustion data often presents itself in tensorial form. Moreover the most valuable datases, i.e. the ones that most closely represent the real physics of the system, are often the ones with the highest dimensionality both in spatial component and in the number of variables. This has sparked interest in methods that process data directly into this form and don't require reshaping into matrix form as in the case of SVD based methods and DMD.

The current report focuses on describing the application of these techniques to enhance the algorithms that allows patter identification in combustion data. High Order Singular Value Decomposition (CITAZIONE), a generalization of SVD, is discussed as a tool to both compress and extract physical patterns from huge combustion datasets. In particular, this method allows for multi modal analysis along each axis of the tensor at the same time, allowing to quantify the importance of each dimension of the tensor.  
When combined with HODMD this, allows for identification of key feature in the dynamics of the reactive flow.

Furthermore, 

%eseguendo una riduzione della dimensionalità indipendente lungo ciascun modo del tensore — dimensioni spaziali, temporali e delle variabili — consentendo migliori tassi di compressione e un filtraggio del rumore più mirato. Combinata con l'HODMD, l'HOSVD fornisce un framework multi-dimensionale naturale per analizzare le dinamiche complesse e multi-scala dei flussi reattivi. Per arricchire ulteriormente l'interpretazione fisica di queste decomposizioni, la Decomposizione di Koopman Spazio-Temporale (STKD) [Le Clainche et al., 2020] estende l'HODMD al dominio spaziale, identificando pattern di onde viaggianti caratterizzati sia da frequenze che da numeri d'onda. Ciò consente di rappresentare il flusso come sovrapposizione di strutture propaganti, offrendo una comprensione più profonda dell'interazione tra dinamiche spaziali e temporali nella combustione turbolenta. Al di là degli approcci basati sulla decomposizione, modelli surrogati data-driven come la Regressione con Processi Gaussiani (GPR) [Rasmussen & Williams, 2006] offrono una prospettiva complementare, consentendo la ricostruzione probabilistica e la predizione dei campi di flusso a partire da osservazioni sparse. Quando estesa con la riduzione della dimensionalità basata su HOSVD, la GPR può operare in uno spazio tensoriale compresso, rendendola scalabile agli dataset ad alta dimensionalità generati da DNS e LES, preservando al contempo le capacità di quantificazione dell'incertezza. In questo lavoro, applichiamo ed estendiamo queste metodologie — HOSVD, HODMD, GPR potenziata con HOSVD e STKD — a dataset di combustione DNS e LES, con l'obiettivo di fornire un framework efficiente, fisicamente interpretabile e computazionalmente trattabile per l'analisi della dinamica dei flussi reattivi.

\end{comment}